\section{The Median-of-Ratios MoM-IV Estimator}
\label{sec:mor}

\begin{algorithm}[H]
\label{alg:mor}
\caption{Median-of-Ratios Estimator}
\KwIn{Sample $(Y_i,X_i,Z_i)_{i=1}^{n}$, number of blocks $k$}
\KwOut{Estimate $\hat{\beta}_{\mathrm{RM}}$}

Partition $(Y_i,X_i,Z_i)_{i=1}^{n}$ into $k$ blocks $B_1, \dots, B_k$ of size $m = \lfloor n/k \rfloor$\;

\For{$j = 1, \dots, k$}{
    Compute the block means $\hat{S}_{ZY}^{(j)} = \frac{1}{m}\sum_{i \in B_j} Z_iY_i$ and $\hat{S}_{ZX}^{(j)} = \frac{1}{m}\sum_{i \in B_j} Z_iX_i$\;
    Compute block level IV estimates $\hat{\beta}^{(j)}=\frac{\hat{S}_{ZY}^{(j)}}{\hat{S}_{ZX}^{(j)}}$\;
}

\Return{$\hat{\beta}_{\mathrm{MR}} = \mathrm{median}(\hat{\beta}^{(1)},\dots,\hat{\beta}^{(k)})$}\;
\end{algorithm}

\begin{lemma}[Block-level error decomposition]
\label{lem:mor_decomp}
Under \eqref{eq:structural} and \ref{ass:exog}, for each block $j$,
\begin{equation}
\label{eq:block_error}
    \hat{\beta}^{(j)} - \beta = \frac{\bar{\varepsilon}_j^Z}{\hat{S}_{ZX}^{(j)}},
\end{equation}
where $\bar{\varepsilon}_j^Z = \frac{1}{m}\sum_{i \in B_j}Z_i\varepsilon_i$ satisfies $\E[\bar{\varepsilon}_j^Z] = 0$.
\end{lemma}

\begin{proof}
Multiply $Y_i = \beta X_i + \varepsilon_i$ by $Z_i$ and average over block $B_j$:
\[
    \hat{S}_{ZY}^{(j)} = \frac{1}{m}\sum_{i \in B_j}Z_i(\beta X_i + \varepsilon_i) = \beta\hat{S}_{ZX}^{(j)} + \bar{\varepsilon}_j^Z.
\]
Dividing by $\hat{S}_{ZX}^{(j)}$ gives $\hat{\beta}^{(j)} = \beta + \bar{\varepsilon}_j^Z/\hat{S}_{ZX}^{(j)}$. By \ref{ass:exog} and linearity of expectation, $\E[\bar{\varepsilon}_j^Z] = \frac{1}{m}\sum_{i \in B_j}\E[Z_i\varepsilon_i] = 0$.
\end{proof}

\begin{theorem}[MoR concentration]
\label{thm:mor}
Assume \ref{ass:exog}--\ref{ass:moments}. Let $\delta \in (0,1)$, $k = \lceil 8\ln(1/\delta) \rceil$, and $m = \lfloor n/k \rfloor$. If the instrument strength condition
\begin{equation}
\label{eq:mor_strength}
    m \geq \frac{32\sigma_{ZX}^2}{\mu_{ZX}^2}
\end{equation}
holds, then
\begin{equation}
\label{eq:mor_bound}
    \Prob\!\left[\abs{\hat{\beta}_{\mathrm{MR}} - \beta} > \frac{4\sqrt{2}\,\sigma_{Z\varepsilon}}{\abs{\mu_{ZX}}\sqrt{m}}\right] \leq \delta.
\end{equation}
\end{theorem}

\begin{proof}
The proof applies the MoM counting argument directly to the block-level IV estimates $\hat{\beta}^{(j)}$.

\medskip\noindent\textbf{Step 1 (Two-event decomposition).} By Lemma~\ref{lem:mor_decomp}, $\abs{\hat{\beta}^{(j)} - \beta} = \abs{\bar{\varepsilon}_j^Z}/\abs{\hat{S}_{ZX}^{(j)}}$. To bound this ratio, we control the numerator above and the denominator below. For a threshold $t > 0$ to be determined, define
\[
    A_j(D) = \left\{\abs{\hat{S}_{ZX}^{(j)} - \mu_{ZX}} \leq D\right\}, \qquad B_j(D,T) = \left\{\abs{\bar{\varepsilon}_Z^{(j)}} \leq tD\right\}.
\]
On $A_j(D)\cap B_j(D,t)$, the denominator satisfies $\abs{\bar{S}_{ZX}}\geq D$ by definition of $A_j(D)$, so
\[
    \abs{\hat{\beta}_{\mathrm{IV}} - \beta}= \frac{\abs{\bar{\varepsilon}_Z^{(j)}}}{\abs{\bar{S}_{ZX}}} \leq \frac{tD}{D} = t.
\]
Hence $\Prob\!\left[\abs{\hat\beta_{MR}-\beta}>t\right] \leq \Prob[A_j(D)^c] + \Prob[B_j(D,t)^c]$ by the union bound. For the remainder of this proof we set $D = \abs{\mu_{ZX}}/2$, the symmetric choice that yields the cleanest constants.

    \medskip\noindent\textbf{Step 2 (Chebyshev bounds).}
    We bound each failure probability in turn.

    The failure event $A_j^c = \{\abs{\bar{S}_{ZX}^{(j)}} < \abs{\mu_{ZX}}/2\}$ implies a large deviation of $\bar{S}_{ZX}^{(j)}$ from its mean. Specifically, if $\abs{\bar{S}_{ZX}^{(j)}} < \abs{\mu_{ZX}}/2$ then, by the reverse triangle inequality,
    \[
        \abs{\bar{S}_{ZX}^{(j)} - \mu_{ZX}} \geq \abs{\mu_{ZX}} - \abs{\bar{S}_{ZX}^{(j)}} > \abs{\mu_{ZX}} - \frac{\abs{\mu_{ZX}}}{2} = \frac{\abs{\mu_{ZX}}}{2}.
    \]
    Hence $A_j^c \subseteq \{\abs{\bar{S}_{ZX}^{(j)}-\mu_{ZX}} > \abs{\mu_{ZX}}/2\}$. Since $\Var(\bar{S}_{ZX}^{(j)}) = \sigma_{ZX}^2/m$, Chebyshev's inequality applied with threshold $\abs{\mu_{ZX}}/2$ gives
    \[
        \Prob[A_j^c] \leq \frac{\sigma_{ZX}^2/m}{\mu_{ZX}^2/4} = \frac{4\sigma_{ZX}^2}{m\mu_{ZX}^2}.
    \]

    The failure event is $B_j^c = \{\abs{\bar{\varepsilon}_Z^{(j)}} > t\abs{\mu_{ZX}}/2\}$. Since $\E[\bar{\varepsilon}_Z^{(j)}] = 0$ and $\Var(\bar{\varepsilon_Z^{(j)}}) = \sigma_{Z\varepsilon}^2/m$, Chebyshev's inequality gives
    \[
        \Prob[B_j^c] \leq \frac{\sigma_{Z\varepsilon}^2/m}{t^2\mu_{ZX}^2/4} = \frac{4\sigma_{Z\varepsilon}^2}{mt^2\mu_{ZX}^2}.
    \]

    The instrument strength condition \eqref{eq:mor_strength} ensures $\Prob[A_j^c] \leq 1/8$. Setting $t = 4\sqrt{2}\,\sigma_{Z\varepsilon}/(|\mu_{ZX}|\sqrt{m})$ gives $\Prob[B_j^c] \leq 1/8$.


    \medskip\noindent\textbf{Step 3 (Allocate the error budget and solve for $t$).}
    We require the total failure probability $\Prob[A^c]+\Prob[B^c]$ to be at most $\frac{1}{4}$, and we allocate half the budget to each event.

    \emph{Budget for $A_j^c$.}
    Setting $\Prob[A^c]\leq 1/8$ requires
    \[
        \frac{4\sigma_{ZX}^2}{m\mu_{ZX}^2} \leq \frac{1}{8}
        \iff
        m \geq \frac{32\sigma_{ZX}^2}{\mu_{ZX}^2},
    \]
    which is exactly condition~\eqref{eq:mor_strength}.

    Setting $\Prob[B_j^c]\leq 1/8$ and solving for $t$:
    \[
        \frac{4\sigma_{Z\varepsilon}^2}{nt^2\mu_{ZX}^2} \leq \frac{1}{8} \implies t \geq \frac{\sqrt{32}\sigma_{Z\varepsilon}}{\abs{\mu_{ZX}}\sqrt{m}}.
    \]
    By the union bound,
    \[
        \Prob\!\left[\abs{\hat{\beta}_{\mathrm{IV}}-\beta} > t\right] \leq \Prob[A^c] + \Prob[B^c] \leq \frac{1}{4} + \frac{1}{4} = \frac{1}{8}.
    \]

\medskip\noindent\textbf{Step 4 (Counting argument).} Define $\zeta_j = \mathbf{1}\{|\hat{\beta}^{(j)} - \beta| > t\}$. By Step~2, $\E[\zeta_j] \leq 1/4$. Since the blocks are disjoint, the $(\zeta_j)_{j=1}^{k}$ are independent. If $|\hat{\beta}_{\mathrm{MR}} - \beta| > t$, the median of the block estimates lies outside $(\beta - t, \beta + t)$, which requires at least $k/2$ blocks to satisfy $|\hat{\beta}^{(j)} - \beta| > t$. Therefore,
\[
    \Prob\!\left[\abs{\hat{\beta}_{\mathrm{MR}} - \beta} > t\right] \leq \Prob\!\left[\sum_{j=1}^{k}\zeta_j \geq \frac{k}{2}\right].
\]

\medskip\noindent\textbf{Step 5 (Hoeffding).} By Hoeffding's inequality, with overshoot $k/2 - k/4 = k/4$,
\[
    \Prob\!\left[\sum_{j=1}^{k}\zeta_j \geq \frac{k}{2}\right] \leq \exp\!\left(\frac{-2(k/4)^2}{k}\right) = e^{-k/8}.
\]
Setting $k = \lceil 8\ln(1/\delta)\rceil$ gives $e^{-k/8} \leq \delta$, completing the proof.
\end{proof}

\begin{remark}[Logarithmic dependence on $\delta$]
\label{rem:iv_polynomial}
    The bound~\eqref{eq:mor_bound} contains the factor $\ln{(1/\delta)}$, which is polynomial in $1/\delta$.
\end{remark}

\begin{remark}[The symmetric threshold is not optimal]
\label{rem:iv_D_choice}
    The choice $D = \abs{\mu_{ZX}}/2$ was made for simplicity. It splits the distance between $0$ and $\abs{\mu_{ZX}}$ evenly, but this is not the split that minimises the error bound for given $n$ and $\delta$. Section~\ref{sec:mor_optimised} shows that treating $D$ as a free parameter and optimising it reduces the instrument strength constant from something to something else and improves the error constant by a factor
    of another thing.
\end{remark}
\todo{add the stuff}

\todo{keep going}


\subsection{Optimised Threshold}
\label{sec:mor_optimised}

The proof of Theorem~\ref{thm:mor} introduced a general threshold $D$ but then
fixed it symetrically. Here we keep $D$ free and choose it to minimise the
resulting error bound symmetrically, at the cost of a more involved Step~3.

\begin{theorem}[Standard MoR concentration, optimised threshold]
\label{thm:mor_opt}
    Assume \ref{ass:exog}--\ref{ass:moments}. Define
    \[
        \rho = \frac{4\sigma_{ZX}^2}{m\mu_{ZX}^2}.
    \]
    For any $\delta\in(0,1)$, if $\rho < 1$ (i.e., $n > \sigma_{ZX}^2/(\delta\mu_{ZX}^2)$),
    then the choice $D^* = (1-\rho^{1/3})\abs{\mu_{ZX}}$ is optimal and
    \begin{equation}
    \label{eq:iv_bound_opt}
        \Prob\!\left[\abs{\hat{\beta}_{\mathrm{IV}} - \beta} > \frac{\sigma_{Z\varepsilon}}{\abs{\mu_{ZX}}\sqrt{\delta n}(1-\rho^{1/3})^{3/2}} \right] \leq \delta.
    \end{equation}
\end{theorem}

\begin{proof}
    Steps~1 and~2 are identical to those of Theorem~\ref{thm:iv} but with $D$ left as a free parameter; we arrive at the Chebyshev bounds
    \[
        \Prob[A(D)^c] \leq \frac{\sigma_{ZX}^2/n}{(\abs{\mu_{ZX}}-D)^2}, \qquad \Prob[B(D,t)^c] \leq \frac{\sigma_{Z\varepsilon}^2/n}{t^2 D^2}.
    \]
    Requiring their sum to be at most $\delta$ gives the joint constraint
    \begin{equation}
    \label{eq:iv_opt_constraint}
        \frac{\sigma_{ZX}^2}{n(\abs{\mu_{ZX}}-D)^2} + \frac{\sigma_{Z\varepsilon}^2}{n t^2 D^2} \leq \delta.
    \end{equation}

    \medskip\noindent\textbf{Step 3 (Optimise $D$, then solve for $t$).}
    Write $D = \eta\abs{\mu_{ZX}}$ with $\eta\in(0,1)$. Dividing~\eqref{eq:iv_opt_constraint} through by $\delta$ and substituting $\rho = \sigma_{ZX}^2/(\delta n\mu_{ZX}^2)$:
    \[
        \frac{\rho}{(1-\eta)^2} + \frac{\sigma_{Z\varepsilon}^2}{\delta n\mu_{ZX}^2\eta^2 t^2} \leq 1.
    \]
    Solving for $t^2$:
    \[
        t^2 \geq
        \frac{\sigma_{Z\varepsilon}^2}{\delta n\mu_{ZX}^2}
        \cdot \frac{1}{f(\eta)},
        \qquad
        f(\eta) = \eta^2\!\left(1 - \frac{\rho}{(1-\eta)^2}\right).
    \]
    Minimising $t$ is equivalent to maximising $f$ over $\eta\in(0,1)$. Differentiating and setting $f'(\eta)=0$:
    \[
        f'(\eta) = 2\eta\!\left(1 - \frac{\rho}{(1-\eta)^2}\right) - \frac{2\rho\eta^2}{(1-\eta)^3} = 0.
    \]
    Dividing through by $2\eta > 0$:
    \[
        1 - \frac{\rho}{(1-\eta)^2} - \frac{\rho\eta}{(1-\eta)^3} = 0 \iff (1-\eta)^3 - \rho(1-\eta) - \rho\eta = 0 \iff (1-\eta)^3 = \rho.
    \]
    Hence $\eta^* = 1 - \rho^{1/3}$, which lies in $(0,1)$ if and only if $\rho < 1$, the stated instrument strength condition. Substituting $(1-\eta^*)^2 = \rho^{2/3}$:
    \[
        f(\eta^*) = (1-\rho^{1/3})^2\!\left(1 - \frac{\rho}{\rho^{2/3}}\right) = (1-\rho^{1/3})^2(1 - \rho^{1/3}) = (1-\rho^{1/3})^3.
    \]
    The minimised threshold is therefore
    \[
        t^* = \frac{\sigma_{Z\varepsilon}}{\abs{\mu_{ZX}}\sqrt{\delta n}(1-\rho^{1/3})^{3/2}},
    \]
    and the union bound gives $\Prob[\abs{\hat\beta_{IV}-\beta}>t^*]\leq\delta$.
\end{proof}

\begin{remark}[Comparison with the symmetric bound]
\label{rem:iv_opt_compare}
    As $\rho\to 0$ (strong instruments or large $n$), the correction factor $(1-\rho^{1/3})^{-3/2}\to 1$ and~\eqref{eq:iv_bound_opt} approaches $\sigma_{Z\varepsilon}/(\abs{\mu_{ZX}}\sqrt{\delta n})$, a factor of $2\sqrt{2}\approx 2.83$ tighter than~\eqref{eq:iv_bound}. As $\rho\to 1$ (approaching the boundary of instrument relevance), the factor diverges, reflecting the genuine difficulty of ratio estimation near instrument irrelevance. The instrument strength condition is eight times weaker than Theorem~\ref{thm:iv}: $\rho < 1$ versus $\rho < 1/8$.
\end{remark}


\subsection{Cantelli Improvement}
\label{sec:iv_cantelli}

The Chebyshev bound applied to $A(D)^c$ in the proofs above is two-sided. However, $A(D)^c = \{\abs{\bar{S}_{ZX}} < D\}$ is a one-sided failure event: only a \emph{downward} deviation of $\bar{S}_{ZX}$ from $\mu_{ZX}$ causes the denominator to be too small; an upward deviation makes it larger and is harmless. Cantelli's inequality exploits this one-sidedness.

\begin{lemma}[Cantelli's inequality]
    \label{lem:cantelli}
    Let $W$ be a random variable with $\E[W]=\mu$ and $\Var(W)=\sigma^2<\infty$.
    For any $\tau>0$,
    \[
        \Prob[W - \mu \leq -\tau] \leq \frac{\sigma^2}{\sigma^2 + \tau^2}.
    \]
    This is strictly smaller than the Chebyshev bound $\sigma^2/\tau^2$ for all $\tau>0$, with the greatest relative gain when $\tau\approx\sigma$.
\end{lemma}

\begin{theorem}[Standard IV concentration, Cantelli improvement]
\label{thm:iv_cantelli}
    Assume \ref{ass:exog}--\ref{ass:moments} and $\mu_{ZX}>0$. Write $\gamma_{\mathrm{IV}} = \sigma_{ZX}^2/(n\mu_{ZX}^2)$ and define
    \begin{equation}
    \label{eq:h_eta}
        h(\eta) = \eta^2
        \frac{\delta\bigl(\gamma_{\mathrm{IV}}+(1-\eta)^2\bigr)-\gamma_{\mathrm{IV}}}
            {\gamma_{\mathrm{IV}}+(1-\eta)^2},
        \quad \eta\in(0,1).
    \end{equation}
    If
    \begin{equation}
    \label{eq:iv_cantelli_strength}
        n > \frac{(1-\delta)\sigma_{ZX}^2}{\delta\mu_{ZX}^2}, \quad\text{equivalently,}\quad \gamma_{\mathrm{IV}} < \frac{\delta}{1-\delta},
    \end{equation}
    then $h(\eta)>0$ on a non-empty interval and the estimator satisfies
    \begin{equation}
    \label{eq:iv_cantelli_bound}
        \Prob\!\left[\abs{\hat{\beta}_{\mathrm{IV}} - \beta}
            > \frac{\sigma_{Z\varepsilon}}{\abs{\mu_{ZX}}\sqrt{nh(\eta^*_{\mathrm{IV}})}}
            \right] \leq \delta,
    \end{equation}
    where $\eta^*_{\mathrm{IV}}$ is the maximiser of $h$ over its feasible domain.
\end{theorem}

\begin{proof}
    We use the same two-event framework with $D = \eta\mu_{ZX}$ kept general.

    \medskip\noindent\textbf{Step 1 (Two-event decomposition).}
    As in Theorem~\ref{thm:iv}, define $A(D) = \{\bar{S}_{ZX}\geq D\}$ and $B(D,t) = \{\abs{\bar{S}_{Z\varepsilon}}\leq tD\}$.
    On $A(D)\cap B(D,t)$, $\abs{\hat\beta_{IV}-\beta}\leq t$.

    \medskip\noindent\textbf{Step 2 (Cantelli on $A^c$, Chebyshev on $B^c$).}

    Since $\mu_{ZX}>0$ and $D<\mu_{ZX}$, the failure event satisfies
    \[
        A(D)^c = \{\bar{S}_{ZX} < D\}
        \subseteq
        \{\bar{S}_{ZX} - \mu_{ZX} < -(\mu_{ZX}-D)\}.
    \]
    This is a one-sided downward deviation, so Lemma~\ref{lem:cantelli} applies with $W=\bar{S}_{ZX}$, $\sigma^2=\sigma_{ZX}^2/n$, and $\tau=\mu_{ZX}-D$:
    \[
        \Prob[A(D)^c]
        \leq
        \frac{\sigma_{ZX}^2/n}{\sigma_{ZX}^2/n + (\mu_{ZX}-D)^2}.
    \]

    The event $B(D,t)^c = \{\abs{\bar{S}_{Z\varepsilon}}>tD\}$ is two-sided (a large residual of either sign causes failure), so Chebyshev is appropriate:
    \[
        \Prob[B(D,t)^c] \leq \frac{\sigma_{Z\varepsilon}^2/n}{t^2 D^2}.
    \]

    \medskip\noindent\textbf{Step 3 (Constraint and optimization).}
    Setting the total failure probability to at most $\delta$ and substituting $D=\eta\mu_{ZX}$, $\gamma_{\mathrm{IV}}=\sigma_{ZX}^2/(n\mu_{ZX}^2)$:
    \[
        \frac{\gamma_{\mathrm{IV}}}{\gamma_{\mathrm{IV}}+(1-\eta)^2}
        + \frac{\sigma_{Z\varepsilon}^2}{n t^2\eta^2\mu_{ZX}^2}
        \leq \delta.
    \]
    Isolating the second term and solving for $t^2$:
    \[
        t^2 \geq \frac{\sigma_{Z\varepsilon}^2}{n\mu_{ZX}^2h(\eta)},
    \]
    where $h(\eta)$ is defined in~\eqref{eq:h_eta}. The function $h(\eta)$ is positive only when $\delta(\gamma_{\mathrm{IV}}+(1-\eta)^2)>\gamma_{\mathrm{IV}}$, i.e., $(1-\eta)^2>\gamma_{\mathrm{IV}}(1-\delta)/\delta$. The feasible region is non-empty if and only if $\gamma_{\mathrm{IV}}<\delta/(1-\delta)$, which is condition~\eqref{eq:iv_cantelli_strength}. Minimizing $t$ over the feasible region by maximizing $h$ and applying the union bound gives~\eqref{eq:iv_cantelli_bound}.
\end{proof}

\begin{remark}[When Cantelli helps]
\label{rem:cantelli_gain}
    Compared with the optimised Chebyshev condition ($n>\sigma_{ZX}^2/(\delta\mu_{ZX}^2)$), the Cantelli condition~\eqref{eq:iv_cantelli_strength} is weaker by a factor of $1/(1-\delta)$. At $\delta=0.05$ this ratio is $1/0.95\approx 1.05$, so the gain is negligible for typical significance levels. The improvement is more meaningful at moderate confidence levels: at $\delta=0.5$ the factor reaches $2$. For the standard IV estimator the Cantelli improvement is therefore primarily of theoretical interest. By contrast, for the MoR estimator (Section~\ref{sec:mor_cantelli}), Cantelli extends the valid range of the instrument strength parameter from $\gamma<1/4$ to $\gamma<1/3$, a qualitatively more substantial gain.
\end{remark}

\begin{remark}[No closed form for $\eta^*_{\mathrm{IV}}$]
    \label{rem:cantelli_noform}
    The first-order condition $h'(\eta)=0$ reduces to a quartic in $\eta$ with no general closed-form solution. The optimised Chebyshev version (Theorem~\ref{thm:iv_opt}) yielded the clean form $\eta^*=1-\rho^{1/3}$ because the Chebyshev bound $\sigma^2/\tau^2$ produces a simpler algebraic structure than the Cantelli bound $\sigma^2/(\sigma^2+\tau^2)$. The bound~\eqref{eq:iv_cantelli_bound} must therefore be evaluated numerically for specific parameter values.
\end{remark}
